\subsection{DEVASTATION OF THE VALLEY}

The war took a cruel new twist in the years 1864 and 1865, when the Union
armies adopted the policy of systematically burning out large areas of the
Confederate "breadbasket" areas (the most famous being Sherman's March to the
Sea, and Sheridan's destruction of the farms in the Shenandoah Valley). This
policy was adopted to deny food supplies to the main Confederate armies, and
also to hinder the operations of Confederate Partisan troops.

\begin{enumerate}
  \item Devasation can only be done by Union Infantry and/or Cavalry units in
        the 1864 scenarios dated August or later (Scenarios \#10 - \#16).
  \item The effects of Devastation cannot be Repaired.
  \item Devastation is handled similarly to Destruction; Movement Allowance
        Factors are expended (use the Destruction column of the MOVEMENT
        AND TERRAIN EFFECTS CHART) for units involved in it. Extra Movement
        Allowance Factors are expended if other Destruction also takes place
        in a hex (i.e. a bridge, rail line, etc.).
  \item Devastation must take place over a wide area in order to be effective.
        Therefore, a "Devastated Area" consists of a hex and all six hexes
        adjacent to it, a total of seven hexes.

        \textbf{Note:} Devastation can be carried out anywhere on the Mapboard.
        However, only certain areas are worth Victory Points to the Union Player.
        Devastated Areas which are worth Victory Points to the Union Player must
        lie between hex rows K to T inclusive, and be over 12 hexes distant from
        north edge of the Mapboard.

        All hexes must have units move into them and expend Movement Allowance
        Factors to Devastate them. The Devastated Area is then marked by placing
        a Destruction Counter in the center hex of the Area.

  \item Each Devastated Area counts towards Victory Points for the Union Player,
        inhibits the Partisans (see Partisan rules), and creates supply problems
        for all units in the area.
\end{enumerate}

END OF ADVANCED GAME RULES